% Statement Bank --- internal supporting obligations.
%
% The public-API obligations of the Jacobian Challenge are stated
% directly in `tex/statements/challenge-statement.tex` (Part 1). The
% blocks below are internal helper statements used in the construction
% of those obligations: the period pairing, the period lattice
% theorem, and the trace-pullback identity on holomorphic 1-forms.
% Their labels are preserved so that downstream sections can continue
% to depend on them.

\section{Statement Bank}
\label{sec:statement-bank}

This bank fixes names for the informal mathematical obligations not
already exposed in the Statement of the Challenge. The labels are chosen
so they can later become Lean-blueprint nodes.

\begin{definition}[Period pairing]
\label{def:period-pairing}
\uses{lem:period-homology-invariance, thm:period-pairing-from-path-integral}
\lean{JacobianChallenge.Periods.periodPairing}
The period pairing is the group homomorphism
\[
  \Pi_X : H_{1}(X,\Z) \longrightarrow \HdualX,\qquad
  \gamma \longmapsto \left(\omega \mapsto \int_{\gamma}\omega\right).
\]
\leanok

\proofplan
\leanok
Define integration functional on cycles, prove path/homology invariance, and
package as an additive homomorphism into \(\HdualX\).
\end{definition}

\begin{layman}
The \emph{period pairing} bundles all the loop integrals at once. Take
any closed loop on the surface (a class in the integer homology --- a
finite-rank free abelian group) and map it to the linear functional
``integrate any holomorphic 1-form along me.'' That gives an element
of the dual of the space of holomorphic 1-forms --- an finite-dimensional linear-algebraic gadget. The pairing sends loops to
such functionals, in a way that respects loop addition (concatenation)
and integer scaling. Why bother? It's the bridge from topology
(homology classes of loops) to analysis (linear functionals on
1-forms). The image of this map will be the period lattice, the
discrete grid that gets quotiented to form the Jacobian.
\end{layman}

\begin{theorem}[Period lattice theorem]
\label{thm:period-lattice}
\uses{input:riemann-bilinear, def:period-pairing, thm:analytic-eq-topological-genus, lem:h1-basis-of-compact-riemann-surface-u, lem:riemann-classical-real-li-input-u}
The image \(\periods_X := \Pi_X(H_1(X,\Z))\) is a full lattice in the
real vector space underlying \(\HdualX\).
\lean{JacobianChallenge.Periods.periodFullComplexLattice}
\leanok
\proofplan
Apply Riemann bilinear relations to show full real rank and discreteness; combine
to obtain a full lattice in the underlying real vector space.
\end{theorem}
\begin{proof}
\leanok
\end{proof}

% Universe-u companions of the period-lattice content (the public Jacobian
% universe bridge consumes these); both are open obligations.
\begin{lemma}[Universe-$u$ singular homology and period-pairing shape]
\label{lem:universe-u-singular-homology-period-pairing-shape}
For a surface \(X : \mathrm{Type}\ u\), the project has a separate
universe-\(u\) singular chain complex, integral \(1\)-cycle object, period
pairing, and basis-aligned period subgroup. These are the shape-level
transport layer used by the universe-\(u\) period lattice route.
\depends{Project-side universe transport in
\code{Jacobian/Blueprint/Sec03/SingularChainComplexZU.lean},
\code{Jacobian/Periods/IntegralOneCycleU.lean},
\code{Jacobian/Periods/PeriodFunctionalU.lean}, and
\code{Jacobian/Periods/BasisAlignedPeriodSubgroupU.lean}; it is grounded in
Mathlib's singular chain and singular homology functors at universe \(u\)
with coefficients \(\mathrm{ULift}\,\Z\).}
\lean{JacobianChallenge.Blueprint.Sec03.singularChainComplexZU,
JacobianChallenge.Blueprint.Sec03.SingularOneChainU,
JacobianChallenge.Blueprint.Sec03.SingularTwoChainU,
JacobianChallenge.Periods.IntegralOneCycleU,
JacobianChallenge.Periods.periodPairing_descent_auxU,
JacobianChallenge.Periods.periodPairingComplexU,
JacobianChallenge.Periods.basisAlignedPeriodSubgroupConcreteU}
\leanok
\end{lemma}
\begin{proof}\leanok\end{proof}

\begin{lemma}[Type-0 H1 basis and symplectic cycle substrate]
\label{lem:type-zero-h1-basis-symplectic-cycle-substrate}
\uses{thm:analytic-eq-topological-genus}
At universe \(0\), the Periods layer supplies a \(\Z\)-basis of integral
first homology indexed by \(2g_{\mathrm{an}}\) and extracts an injective
family of \(2g\) cycles from it. This is the surface-classification/H1
substrate mirrored by the universe-\(u\) H1-basis frontier.
\depends{Project-side Type-0 assembly in
\code{Jacobian/Periods/PeriodFunctional.lean}; it routes through the
surface-classification and analytic-equals-topological-genus bridge already
tracked by \ref{thm:analytic-eq-topological-genus}.}
\lean{JacobianChallenge.Periods.h1_basis_of_compact_riemann_surface,
JacobianChallenge.Periods.symplectic_basis_of_cycles}
\leanok
\end{lemma}
\begin{proof}\leanok\end{proof}

\begin{lemma}[Universe-$u$ $H_1$ basis (open)]
\label{lem:h1-basis-of-compact-riemann-surface-u}
\uses{lem:universe-u-singular-homology-period-pairing-shape, lem:type-zero-h1-basis-symplectic-cycle-substrate}
A \(\Z\)-basis of \(H_1\) for a universe-\(u\) compact Riemann surface
(\(\code{IntegralOneCycleU}\)). This is the universe-\(u\) rerun of the
surface-classification/cellular-homology rank-\(2g\) argument, not another
Riemann-bilinear analytic gap. Open obligation (\code{sorry}).
Tracking: \ghissue{240}.
\lean{JacobianChallenge.Periods.h1_basis_of_compact_riemann_surfaceU}
\end{lemma}

\begin{lemma}[Universe-$u$ Riemann-bilinear $\R$-LI input (open)]
\label{lem:riemann-classical-real-li-input-u}
\uses{lem:riemann-classical-real-li-input, lem:universe-u-singular-homology-period-pairing-shape}
Universe-\(u\) companion of the Riemann-bilinear \(\R\)-linear-independence
input. Its analytic content is the same Riemann-bilinear / Hodge-positivity
frontier as the Type-0 input, now stated for \(\code{IntegralOneCycleU}\) and
\(\code{periodPairingComplexU}\). Open obligation (\code{sorry}).
Tracking: \ghissue{241}.
\lean{JacobianChallenge.Periods.riemann_classical_real_LI_inputU}
\end{lemma}

\begin{layman}
The image of the period pairing --- the set of all integer-coefficient
combinations of period vectors --- is a \emph{full lattice} inside the
real vector space underlying the dual of the holomorphic 1-forms.
``Full'' means it's a discrete grid spanning the whole real space, not
a sublattice of lower rank. This is a direct consequence of the linear
independence of the \(2g\) symplectic-basis period vectors (from the
Riemann bilinear relations), since \(2g\) independent vectors in a
\(2g\)-real-dimensional space generate a lattice of full rank. Why
bother? A full lattice is exactly the input the complex-torus package
needs to produce a compact quotient. With the period lattice in
hand, the Jacobian (Definition~\ref{def:analytic-jacobian}) is just
``dual-of-1-forms modulo lattice'' --- and that quotient is
automatically a compact complex Lie additive group.
\end{layman}
\begin{theorem}[Trace-pullback identity]
\label{thm:trace-pullback}
\uses{def:degree-trace, lem:trace-forms, lem:pullback-forms}
For every \(\eta\in\HY\),
\[
  \tr_f(f^{*}\eta)=\deg(f)\,\eta.
\]
\lean{JacobianChallenge.TraceDegree.analyticPushforward_analyticPullback}
\leanok
\proofplan
Use the fundamental identity \(\tr_f \circ f^* = \deg(f)\) on holomorphic 1-forms
(Theorem~\ref{thm:trace-pullback-identity}) and descend to the Jacobian quotient.
\end{theorem}

\begin{layman}
Take a holomorphic 1-form on the target. Pull it back along the map
to get a 1-form on the source. Trace it back down to the target. The
result equals the original 1-form multiplied by the degree of the
map. The intuition: pulling back creates \(d\) copies of the form (one
per inverse branch in a generic fiber); summing them via the trace
recovers \(d\) times the original. Why bother? This is the algebraic
backbone of the pushforward-pullback identity
(Theorem~\ref{thm:pushforward-pullback}). After dualizing and
descending to Jacobians, it becomes ``push-then-pull equals
multiplication-by-degree'' on Jacobian elements --- forcing push, pull,
and degree to interact through the classical trace identity rather
than through arbitrary homomorphisms invented to game the challenge.
\end{layman}
